% Generated draft: requires mathematical and source review.
\documentclass[11pt]{article}
\usepackage[margin=25mm]{geometry}
\usepackage{fontspec}
\setmainfont{Latin Modern Roman}
\usepackage{amsmath,amssymb,mathtools,unicode-math}
\setmathfont{Latin Modern Math}
\usepackage{enumitem,needspace,xurl,hyperref}
\setlength{\parindent}{0pt}
\setlength{\parskip}{5pt}
\setlength{\emergencystretch}{3em}
\newcommand{\N}{\mathbb{N}}
\newcommand{\Z}{\mathbb{Z}}
\newcommand{\Q}{\mathbb{Q}}
\newcommand{\R}{\mathbb{R}}
\newcommand{\C}{\mathbb{C}}
\newcommand{\F}{\mathbb{F}}
\newcommand{\A}{\mathbb{A}}
\newcommand{\PP}{\mathbb{P}}
\newcommand{\E}{\mathbb{E}}
\newcommand{\eps}{\varepsilon}
\newcommand{\dd}{\,\mathrm{d}}
\begin{document}
\section*{UnsolvedMath ID: 801060}
\textbf{Source problem code:} catalogue entry
\subsection{Definitions and mathematical statement}
Let \(H_{a,b}\) denote the H\'enon family referred to in the supplied source, and let \(\mathcal{WY}\) denote its Wang--Young parameter set. The source describes \(\mathcal{WY}\) as having positive measure and consisting of strongly dissipative maps for which Wang and Young proved the existence of a strange attractor. The formula for \(H_{a,b}\), the parameter domain, the reference measure, the precise dissipativity assumptions, and the definition of strange attractor are not supplied.

The source states that reference [BSt] develops a kneading theory for these maps. For each \((a,b)\in\mathcal{WY}\), write \(\mathfrak K(a,b)\) for the set of kneading sequences assigned to \(H_{a,b}\) by that theory. This notation is a placeholder for the source's construction: its symbolic alphabet, indexing, coding rules, and boundary conventions are not available in the supplied record.

\textbf{Problem.} Describe the set of kneading sequences \(\mathfrak K\) of the H\'enon maps \(H_{a,b}\) with \((a,b)\in\mathcal{WY}\), using the kneading theory of [BSt].

The intended scope of this description remains unresolved. It could ask for a description of \(\mathfrak K(a,b)\) for every parameter pair, or for a characterization of the aggregate realizability set
\[
\mathfrak K_{\mathcal{WY}}
:=\bigcup_{(a,b)\in\mathcal{WY}}\mathfrak K(a,b).
\]
The supplied wording does not explicitly distinguish these targets. A self-contained precise formulation requires the definitions from [BSt] and [WY] and clarification of this scope.
\subsection{Short English statement}
Describe which kneading sequences arise for the H\'enon maps in the Wang--Young parameter set, using the kneading theory cited by the source. The supplied record does not define the coding or parameter set, and leaves unclear whether the requested description is parameter by parameter or across the entire parameter set.
\subsection{Sources}
\begin{itemize}
\item Local UnsolvedMath registry, permanent ID \(801060\), source problem code \texttt{catalogue entry}, title \emph{Henon Maps: Open Problems --- Problem 12}. The record attributes the problem to Berger and Bedford, with proposed year 2023, and labels it open. This is a supplied status label, not an independently verified assessment.
\item Berger and Bedford, \emph{H\'enon Maps: A List of Open Problems} (2023), Problem 12. Supplied source URL: \url{https://arxiv.org/abs/2312.03907}. The record reports access on 2026-10-16 and extraction by \texttt{arxiv\_latex}; these are supplied provenance details, not a claim of independent access or verification.
\item Reference [BSt], cited in the supplied statement for a kneading theory on \(\mathcal{WY}\). No full bibliographic entry or separate URL is supplied.
\item Reference [WY], cited in the supplied statement for the Wang--Young parameter set and existence of a strange attractor. No full bibliographic entry or separate URL is supplied.
\end{itemize}
\end{document}
